\documentclass[12pt, a4paper]{ctexart}
\usepackage{fontspec}
\usepackage{xcolor}
\usepackage{hyperref}
\usepackage{geometry}
\usepackage{enumitem}
\usepackage{parskip}
\usepackage{titlesec}
\usepackage{tcolorbox}
\tcbuselibrary{breakable}
\usepackage{etoolbox}
\usepackage{setspace}
\usepackage{listings}

\geometry{left=2.5cm, right=2.5cm, top=2.5cm, bottom=2.5cm}

\setmainfont{Times New Roman}
\setCJKmainfont{SimSun}

\hypersetup{
    colorlinks=true,
    linkcolor=blue!70!black,
    urlcolor=cyan,
    pdftitle={网络编程讲义},
    pdfauthor={专升本-fifine老师}
}

\definecolor{myblue}{RGB}{30,100,200}
\definecolor{lightblue}{RGB}{230,240,255}
\definecolor{darkblue}{RGB}{0,50,120}

\newtcolorbox{bluebox}[1][]{
    colback=lightblue,
    colframe=myblue,
    coltitle=darkblue,
    fonttitle=\bfseries,
    title={#1},
    boxrule=0.8pt,
    arc=4pt,
    left=6pt,
    right=6pt,
    top=6pt,
    bottom=6pt,
    before skip=8pt,
    after skip=8pt,
    breakable
}

\BeforeBeginEnvironment{bluebox}{\par\vfill}%
\AfterEndEnvironment{bluebox}{\par\vfill}%

\titleformat{\section}
  {\normalfont\Large\bfseries\color{myblue}}{\thesection}{1em}{}
\titlespacing*{\section}{0pt}{3.5ex plus 1ex minus .2ex}{1.5ex plus .2ex}

\title{\textcolor{myblue}{\textbf{网络编程 - 深度解析讲义}}}
\author{专升本-fifine老师 教学组}
\date{2026年03月02日}

\lstset{
    language=Java,
    basicstyle=\ttfamily\small,
    keywordstyle=\color{blue},
    commentstyle=\color{green!60!black},
    stringstyle=\color{orange},
    numbers=left,
    numberstyle=\tiny\color{gray},
    frame=single,
    breaklines=true,
    columns=fullflexible,
    showstringspaces=false
}

\newcommand{\code}[1]{\texttt{#1}}

\begin{document}

\maketitle
\thispagestyle{empty}
\newpage

\section{网络编程}

\begin{enumerate}[label=\textbf{\arabic*.}]

	\item \textbf{以下关于Java中Socket编程的说法错误的是}
	      \begin{enumerate}[label=\Alph*.]
		      \item 用于网络通信
		      \item ServerSocket用于服务器端
		      \item Socket用于客户端
		      \item 不需要处理异常
	      \end{enumerate}

	      \begin{bluebox}{答案与深度解析}
		      \textbf{答案：D}

		      % ======== 阶段一：核心认知框架 ========
		      \textbf{【核心认知框架】}
		      \begin{itemize}[label={}, leftmargin=*, nosep]
			      \item \textbf{题目本质}：考查Socket网络编程的基本概念
			      \item \textbf{关键区分点}：网络编程必须处理异常
			      \item \textbf{命题人意图}：测试对网络编程基础的理解
		      \end{itemize}

		      % ======== 阶段二：选项深度拆解 ========
		      \textbf{【选项原理深度拆解】}
					  \begin{itemize}[label={}, leftmargin=*, nosep]

		      % ---------- 选项A ----------
		      \item \textbf{A. 用于网络通信 — 正确}
		      \begin{itemize}[label={}, leftmargin=*, nosep]
			      \item \textbf{Socket}：网络通信的端点
			      \item \textbf{作用}：建立TCP/UDP连接，进行数据交换
		      \end{itemize}

		      % ---------- 选项B ----------
		      \item \textbf{B. ServerSocket用于服务器端 — 正确}
		      \begin{itemize}[label={}, leftmargin=*, nosep]
			      \item \textbf{ServerSocket}：服务器端监听端口
			            \begin{lstlisting}[language=Java]
ServerSocket serverSocket = new ServerSocket(8080);
Socket client = serverSocket.accept(); // 等待客户端连接
      \end{lstlisting}
		      \end{itemize}

		      % ---------- 选项C ----------
		      \item \textbf{C. Socket用于客户端 — 正确}
		      \begin{itemize}[label={}, leftmargin=*, nosep]
			      \item \textbf{Socket}：客户端连接服务器
			            \begin{lstlisting}[language=Java]
Socket socket = new Socket("localhost", 8080);
      \end{lstlisting}
		      \end{itemize}

		      % ---------- 选项D ----------
		      \item \textbf{D. 不需要处理异常 — 错误}
		      \begin{itemize}[label={}, leftmargin=*, nosep]
			      \item \textbf{必须处理异常}：
			            \begin{itemize}[label={}, nosep]
				            \item IOException：网络通信异常
				            \item SocketException：Socket相关异常
				            \item UnknownHostException：未知主机
			            \end{itemize}
			      \item \textbf{正确写法}：
			            \begin{lstlisting}[language=Java]
try (Socket socket = new Socket("localhost", 8080)) {
    // 通信
} catch (IOException e) {
    e.printStackTrace();
}
      \end{lstlisting}
		      \end{itemize}

		      \end{itemize}

					  % ======== 阶段三：应背诵知识点 ========
		      \textbf{【应背诵知识点（命题人思维版）】}
		      \begin{enumerate}[leftmargin=*, nosep]
			      \item \textbf{Socket通信流程}：
			            \begin{itemize}[label={}, nosep]
				            \item 服务器：ServerSocket监听 → accept() → 获得Socket
				            \item 客户端：new Socket() → 连接服务器
			            \end{itemize}

			      \item \textbf{高频陷阱清单}：
			            \begin{itemize}[label={}, nosep]
				            \item 陷阱1：网络编程必须捕获IOException
			            \end{itemize}
		      \end{enumerate}

		      % ======== 阶段四：命题趋势与实战策略 ========
		      \textbf{【命题趋势与实战策略】}

		      \textbf{考场秒杀技巧}：Socket编程异常处理是必须的
	      \end{bluebox}

	\item \textbf{以下哪个方法用于发送数据？}
	      \begin{enumerate}[label=\Alph*.]
		      \item send()
		      \item write()
		      \item output()
		      \item transmit()
	      \end{enumerate}

	      \begin{bluebox}{答案与深度解析}
		      \textbf{答案：B}

		      % ======== 阶段一：核心认知框架 ========
		      \textbf{【核心认知框架】}
		      \begin{itemize}[label={}, leftmargin=*, nosep]
			      \item \textbf{题目本质}：考查Socket数据发送方法
			      \item \textbf{关键区分点}：使用OutputStream的write方法
			      \item \textbf{命题人意图}：测试对Socket I/O API的掌握
		      \end{itemize}

		      % ======== 阶段二：选项深度拆解 ========
		      \textbf{【选项原理深度拆解】}
					  \begin{itemize}[label={}, leftmargin=*, nosep]

		      % ---------- 选项A ----------
		      \item \textbf{A. send() — 错误}
		      \begin{itemize}[label={}, leftmargin=*, nosep]
			      \item \textbf{辨析}：DatagramPacket有send()，但Socket用Stream
		      \end{itemize}

		      % ---------- 选项B ----------
		      \item \textbf{B. write() — 正确}
		      \begin{itemize}[label={}, leftmargin=*, nosep]
			      \item \textbf{使用OutputStream}：
			            \begin{lstlisting}[language=Java]
Socket socket = new Socket("localhost", 8080);
OutputStream out = socket.getOutputStream();
out.write("Hello".getBytes()); // 发送数据
out.flush(); // 刷新缓冲区
      \end{lstlisting}
		      \end{itemize}

		      % ---------- 选项C ----------
		      \item \textbf{C. output() — 错误}
		      \begin{itemize}[label={}, leftmargin=*, nosep]
			      \item \textbf{辨析}：不是Socket的方法
		      \end{itemize}

		      % ---------- 选项D ----------
		      \item \textbf{D. transmit() — 错误}
		      \begin{itemize}[label={}, leftmargin=*, nosep]
			      \item \textbf{辨析}：不存在此方法
		      \end{itemize}

		      \end{itemize}

					  % ======== 阶段三：应背诵知识点 ========
		      \textbf{【应背诵知识点（命题人思维版）】}
		      \begin{enumerate}[leftmargin=*, nosep]
			      \item \textbf{Socket I/O方法}：
			            \begin{itemize}[label={}, nosep]
				            \item getOutputStream()：获取输出流（发送）
				            \item getInputStream()：获取输入流（接收）
			            \end{itemize}

			      \item \textbf{注意}：write后要flush确保数据发出
		      \end{enumerate}

		      % ======== 阶段四：命题趋势与实战策略 ========
		      \textbf{【命题趋势与实战策略】}

		      \textbf{考场秒杀技巧}：Socket发送用OutputStream.write()
	      \end{bluebox}

	\item \textbf{以下关于Java中URL类的说法错误的是}
	      \begin{enumerate}[label=\Alph*.]
		      \item 用于表示统一资源定位符
		      \item 可以获取网络资源
		      \item 不能解析URL
		      \item 可以获取URL的各个部分
	      \end{enumerate}

	      \begin{bluebox}{答案与深度解析}
		      \textbf{答案：C}

		      % ======== 阶段一：核心认知框架 ========
		      \textbf{【核心认知框架】}
		      \begin{itemize}[label={}, leftmargin=*, nosep]
			      \item \textbf{题目本质}：考查java.net.URL类的功能
			      \item \textbf{关键区分点}：URL类可以解析URL字符串
			      \item \textbf{命题人意图}：测试对URL类的理解
		      \end{itemize}

		      % ======== 阶段二：选项深度拆解 ========
		      \textbf{【选项原理深度拆解】}
					  \begin{itemize}[label={}, leftmargin=*, nosep]

		      % ---------- 选项A ----------
		      \item \textbf{A. 用于表示统一资源定位符 — 正确}
		      \begin{itemize}[label={}, leftmargin=*, nosep]
			      \item \textbf{URL格式}：protocol://host:port/path
		      \end{itemize}

		      % ---------- 选项B ----------
		      \item \textbf{B. 可以获取网络资源 — 正确}
		      \begin{itemize}[label={}, leftmargin=*, nosep]
			      \item \textbf{openStream()}：打开网络连接获取输入流
			            \begin{lstlisting}[language=Java]
URL url = new URL("https://www.example.com");
InputStream in = url.openStream();
// 读取网络资源
      \end{lstlisting}
		      \end{itemize}

		      % ---------- 选项C ----------
		      \item \textbf{C. 不能解析URL — 错误}
		      \begin{itemize}[label={}, leftmargin=*, nosep]
			      \item \textbf{恰恰相反}：URL类专门用于解析URL
			            \begin{lstlisting}[language=Java]
URL url = new URL("https://localhost:8080/path");
url.getProtocol(); // https
url.getHost(); // localhost
url.getPort(); // 8080
url.getPath(); // /path
      \end{lstlisting}
		      \end{itemize}

		      % ---------- 选项D ----------
		      \item \textbf{D. 可以获取URL的各个部分 — 正确}
		      \begin{itemize}[label={}, leftmargin=*, nosep]
			      \item \textbf{获取方法}：getProtocol()、getHost()、getPort()、getPath()等
		      \end{itemize}

		      \end{itemize}

					  % ======== 阶段三：应背诵知识点 ========
		      \textbf{【应背诵知识点（命题人思维版）】}
		      \begin{enumerate}[leftmargin=*, nosep]
			      \item \textbf{URL类常用方法}：
			            \begin{tabular}{|c|c|}
				            \hline
				            \textbf{方法}   & \textbf{返回} \\
				            \hline
				            getProtocol() & 协议          \\
				            \hline
				            getHost()     & 主机名         \\
				            \hline
				            getPort()     & 端口          \\
				            \hline
				            getPath()     & 路径          \\
				            \hline
			            \end{tabular}
		      \end{enumerate}

		      % ======== 阶段四：命题趋势与实战策略 ========
		      \textbf{【命题趋势与实战策略】}

		      \textbf{考场秒杀技巧}：URL类专门解析和获取URL各部分
	      \end{bluebox}

	\item \textbf{以下哪个方法用于获取URL的协议部分？}
	      \begin{enumerate}[label=\Alph*.]
		      \item getProtocol()
		      \item getScheme()
		      \item getType()
		      \item getProtocolType()
	      \end{enumerate}

	      \begin{bluebox}{答案与深度解析}
		      \textbf{答案：B}

		      % ======== 阶段一：核心认知框架 ========
		      \textbf{【核心认知框架】}
		      \begin{itemize}[label={}, leftmargin=*, nosep]
			      \item \textbf{题目本质}：考查URL类获取协议的方法
			      \item \textbf{关键区分点}：getScheme()是正确方法名
			      \item \textbf{命题人意图}：测试对URL API的熟悉程度
		      \end{itemize}

		      % ======== 阶段二：选项深度拆解 ========
		      \textbf{【选项原理深度拆解】}
					  \begin{itemize}[label={}, leftmargin=*, nosep]

		      % ---------- 选项A ----------
		      \item \textbf{A. getProtocol() — 错误}
		      \begin{itemize}[label={}, leftmargin=*, nosep]
			      \item \textbf{辨析}：不是URL类的方法
		      \end{itemize}

		      \item \textbf{B. getScheme() — 正确}
		      \begin{itemize}[label={}, leftmargin=*, nosep]
			      \item \textbf{返回协议}：http、https、ftp等
			            \begin{lstlisting}[language=Java]
URL url = new URL("https://www.example.com");
String scheme = url.getScheme(); // https
      \end{lstlisting}
		      \end{itemize}

		      % ---------- 选项C ----------
		      \item \textbf{C. getType() — 错误}
		      \begin{itemize}[label={}, leftmargin=*, nosep]
			      \item \textbf{辨析}：不存在此方法
		      \end{itemize}

		      % ---------- 选项D ----------
		      \item \textbf{D. getProtocolType() — 错误}
		      \begin{itemize}[label={}, leftmargin=*, nosep]
			      \item \textbf{辨析}：不存在此方法
		      \end{itemize}

		      \end{itemize}

					  % ======== 阶段三：应背诵知识点 ========
		      \textbf{【应背诵知识点（命题人思维版）】}
		      \begin{enumerate}[leftmargin=*, nosep]
			      \item \textbf{记忆技巧}：URL遵循URI规范，用scheme表示协议
		      \end{enumerate}

		      % ======== 阶段四：命题趋势与实战策略 ========
		      \textbf{【命题趋势与实战策略】}

		      \textbf{考场秒杀技巧}：URL获取协议用getScheme()
	      \end{bluebox}

	\item \textbf{以下关于Java中Servlet的说法错误的是}
	      \begin{enumerate}[label=\Alph*.]
		      \item 运行在服务器端
		      \item 是一个接口
		      \item 用于处理HTTP请求
		      \item 可以直接被实例化
	      \end{enumerate}

	      \begin{bluebox}{答案与深度解析}
		      \textbf{答案：D}

		      % ======== 阶段一：核心认知框架 ========
		      \textbf{【核心认知框架】}
		      \begin{itemize}[label={}, leftmargin=*, nosep]
			      \item \textbf{题目本质}：考查Servlet生命周期管理
			      \item \textbf{关键区分点}：Servlet由容器管理，不直接new
			      \item \textbf{命题人意图}：测试对Servlet生命周期的理解
		      \end{itemize}

		      % ======== 阶段二：选项深度拆解 ========
		      \textbf{【选项原理深度拆解】}
					  \begin{itemize}[label={}, leftmargin=*, nosep]

		      % ---------- 选项A ----------
		      \item \textbf{A. 运行在服务器端 — 正确}
		      \begin{itemize}[label={}, leftmargin=*, nosep]
			      \item \textbf{容器}：Tomcat、Jetty等Web容器
		      \end{itemize}

		      % ---------- 选项B ----------
		      \item \textbf{B. 是一个接口 — 正确}
		      \begin{itemize}[label={}, leftmargin=*, nosep]
			      \item \textbf{javax.servlet.Servlet}：核心接口
		      \end{itemize}

		      % ---------- 选项C ----------
		      \item \textbf{C. 用于处理HTTP请求 — 正确}
		      \begin{itemize}[label={}, leftmargin=*, nosep]
			      \item \textbf{HttpServlet}：处理HTTP请求的抽象类
			            \begin{lstlisting}[language=Java]
public class MyServlet extends HttpServlet {
    protected void doGet(HttpServletRequest req, 
                         HttpServletResponse resp) {
        // 处理GET请求
    }
}
      \end{lstlisting}
		      \end{itemize}

		      % ---------- 选项D ----------
		      \item \textbf{D. 可以直接被实例化 — 错误}
		      \begin{itemize}[label={}, leftmargin=*, nosep]
			      \item \textbf{容器管理}：Servlet由Web容器实例化、初始化、调用
			            \begin{itemize}[label={}, nosep]
				            \item 容器创建Servlet实例
				            \item 调用init()初始化
				            \item 调用service()处理请求
				            \item 调用destroy()销毁
			            \end{itemize}
			      \item \textbf{开发中不直接new}
		      \end{itemize}

		      \end{itemize}

					  % ======== 阶段三：应背诵知识点 ========
		      \textbf{【应背诵知识点（命题人思维版）】}
		      \begin{enumerate}[leftmargin=*, nosep]
			      \item \textbf{Servlet生命周期}：
			            \begin{enumerate}[label={}, leftmargin=*, nosep]
				            \item 加载：容器加载Servlet类
				            \item 实例化：new Servlet()
				            \item 初始化：init()
				            \item 处理请求：service()
				            \item 销毁：destroy()
			            \end{enumerate}

			      \item \textbf{高频陷阱清单}：
			            \begin{itemize}[label={}, nosep]
				            \item 陷阱1：开发者不直接new Servlet，由容器管理
			            \end{itemize}
		      \end{enumerate}

		      % ======== 阶段四：命题趋势与实战策略 ========
		      \textbf{【命题趋势与实战策略】}

		      \textbf{考场秒杀技巧}：Servlet由容器管理生命周期
	      \end{bluebox}

	\item \textbf{以下哪个方法是Servlet的生命周期方法？}
	      \begin{enumerate}[label=\Alph*.]
		      \item init()
		      \item start()
		      \item execute()
		      \item process()
	      \end{enumerate}

	      \begin{bluebox}{答案与深度解析}
		      \textbf{答案：A}

		      % ======== 阶段一：核心认知框架 ========
		      \textbf{【核心认知框架】}
		      \begin{itemize}[label={}, leftmargin=*, nosep]
			      \item \textbf{题目本质}：考查Servlet生命周期方法
			      \item \textbf{关键区分点}：init/service/destroy是生命周期方法
			      \item \textbf{命题人意图}：测试对Servlet生命周期的掌握
		      \end{itemize}

		      % ======== 阶段二：选项深度拆解 ========
		      \textbf{【选项原理深度拆解】}
					  \begin{itemize}[label={}, leftmargin=*, nosep]

		      % ---------- 选项A ----------
		      \item \textbf{A. init() — 正确}
		      \begin{itemize}[label={}, leftmargin=*, nosep]
			      \item \textbf{初始化方法}：容器调用一次
			            \begin{lstlisting}[language=Java]
public void init(ServletConfig config) throws ServletException {
    // 初始化代码
}
      \end{lstlisting}
		      \end{itemize}

		      % ---------- 选项B ----------
		      \item \textbf{B. start() — 错误}
		      \begin{itemize}[label={}, leftmargin=*, nosep]
			      \item \textbf{辨析}：不是Servlet的方法
		      \end{itemize}

		      % ---------- 选项C ----------
		      \item \textbf{C. execute() — 错误}
		      \begin{itemize}[label={}, leftmargin=*, nosep]
			      \item \textbf{辨析}：不是Servlet的方法
		      \end{itemize}

		      % ---------- 选项D ----------
		      \item \textbf{D. process() — 错误}
		      \begin{itemize}[label={}, leftmargin=*, nosep]
			      \item \textbf{辨析}：不是Servlet的方法
		      \end{itemize}

		      \end{itemize}

					  % ======== 阶段三：应背诵知识点 ========
		      \textbf{【应背诵知识点（命题人思维版）】}
		      \begin{enumerate}[leftmargin=*, nosep]
			      \item \textbf{Servlet生命周期方法}：
			            \begin{itemize}[label={}, nosep]
				            \item init()：初始化
				            \item service()：处理请求（doGet/doPost）
				            \item destroy()：销毁
			            \end{itemize}

			      \item \textbf{调用次数}：
			            \begin{itemize}[label={}, nosep]
				            \item init()：一次
				            \item service()：每次请求
				            \item destroy()：一次
			            \end{itemize}
		      \end{enumerate}

		      % ======== 阶段四：命题趋势与实战策略 ========
		      \textbf{【命题趋势与实战策略】}

		      \textbf{考场秒杀技巧}：Servlet生命周期 = init + service + destroy
	      \end{bluebox}

	\item \textbf{以下关于Java中JSP的说法错误的是}
	      \begin{enumerate}[label=\Alph*.]
		      \item 本质是一个Servlet
		      \item 可以嵌入Java代码
		      \item 性能比Servlet高
		      \item 用于生成动态网页
	      \end{enumerate}

	      \begin{bluebox}{答案与深度解析}
		      \textbf{答案：C}

		      % ======== 阶段一：核心认知框架 ========
		      \textbf{【核心认知框架】}
		      \begin{itemize}[label={}, leftmargin=*, nosep]
			      \item \textbf{题目本质}：考查JSP与Servlet的关系
			      \item \textbf{关键区分点}：JSP首次访问需编译为Servlet
			      \item \textbf{命题人意图}：测试对JSP机制的理解
		      \end{itemize}

		      % ======== 阶段二：选项深度拆解 ========
		      \textbf{【选项原理深度拆解】}
					  \begin{itemize}[label={}, leftmargin=*, nosep]

		      % ---------- 选项A ----------
		      \item \textbf{A. 本质是一个Servlet — 正确}
		      \begin{itemize}[label={}, leftmargin=*, nosep]
			      \item \textbf{编译过程}：JSP → Servlet → 编译 → class
			            \begin{itemize}[label={}, nosep]
				            \item 首次访问时，容器将JSP翻译成Java Servlet源码
				            \item 然后编译成.class文件
			            \end{itemize}
		      \end{itemize}

		      % ---------- 选项B ----------
		      \item \textbf{B. 可以嵌入Java代码 — 正确}
		      \begin{itemize}[label={}, leftmargin=*, nosep]
			      \item \textbf{JSP脚本元素}：
			            \begin{itemize}[label={}, nosep]
				            \item <\% \%>：代码片段
				            \item <\%= \%>：表达式输出
				            \item <\%! \%>：声明
			            \end{itemize}
		      \end{itemize}

		      % ---------- 选项C ----------
		      \item \textbf{C. 性能比Servlet高 — 错误}
		      \begin{itemize}[label={}, leftmargin=*, nosep]
			      \item \textbf{性能分析}：
			            \begin{itemize}[label={}, nosep]
				            \item 首次访问：JSP需要编译，性能低于直接使用Servlet
				            \item 后续访问：编译后与Servlet性能相当
				            \item 总体：Servlet性能更稳定
			            \end{itemize}
		      \end{itemize}

		      % ---------- 选项D ----------
		      \item \textbf{D. 用于生成动态网页 — 正确}
		      \begin{itemize}[label={}, leftmargin=*, nosep]
			      \item \textbf{应用场景}：视图层，动态生成HTML
		      \end{itemize}

		      \end{itemize}

					  % ======== 阶段三：应背诵知识点 ========
		      \textbf{【应背诵知识点（命题人思维版）】}
		      \begin{enumerate}[leftmargin=*, nosep]
			      \item \textbf{JSP vs Servlet}：
			            \begin{tabular}{|c|c|c|}
				            \hline
				            \textbf{对比} & \textbf{JSP} & \textbf{Servlet} \\
				            \hline
				            角色          & 视图层          & 控制层              \\
				            \hline
				            编写难度        & 简单           & 较难               \\
				            \hline
				            首次性能        & 差            & 好                \\
				            \hline
			            \end{tabular}

			      \item \textbf{高频陷阱清单}：
			            \begin{itemize}[label={}, nosep]
				            \item 陷阱1：JSP首次访问需编译，首次性能差
			            \end{itemize}
		      \end{enumerate}

		      % ======== 阶段四：命题趋势与实战策略 ========
		      \textbf{【命题趋势与实战策略】}

		      \textbf{考场秒杀技巧}：JSP首次访问编译，Servlet性能更稳定
	      \end{bluebox}

	\item \textbf{以下哪个是JSP中的指令？}
	      \begin{enumerate}[label=\Alph*.]
		      \item <jsp:include>
		      \item <jsp:forward>
		      \item <\%@ \%>
		      \item <jsp:param>
	      \end{enumerate}

	      \begin{bluebox}{答案与深度解析}
		      \textbf{答案：C}

		      % ======== 阶段一：核心认知框架 ========
		      \textbf{【核心认知框架】}
		      \begin{itemize}[label={}, leftmargin=*, nosep]
			      \item \textbf{题目本质}：考查JSP指令与动作的区别
			      \item \textbf{关键区分点}：指令使用<\%@ \%>语法
			      \item \textbf{命题人意图}：测试对JSP语法的掌握
		      \end{itemize}

		      % ======== 阶段二：选项深度拆解 ========
		      \textbf{【选项原理深度拆解】}
					  \begin{itemize}[label={}, leftmargin=*, nosep]

		      % ---------- 选项A ----------
		      \item \textbf{A. <jsp:include> — JSP动作}
		      \begin{itemize}[label={}, leftmargin=*, nosep]
			      \item \textbf{类型}：JSP Action（动作）
			      \item \textbf{作用}：动态包含
			            \begin{lstlisting}[language=Java]
<jsp:include page="header.jsp" />
      \end{lstlisting}
		      \end{itemize}

		      % ---------- 选项B ----------
		      \item \textbf{B. <jsp:forward> — JSP动作}
		      \begin{itemize}[label={}, leftmargin=*, nosep]
			      \item \textbf{类型}：JSP Action
			      \item \textbf{作用}：请求转发
			            \begin{lstlisting}[language=Java]
<jsp:forward page="next.jsp" />
      \end{lstlisting}
		      \end{itemize}

		      % ---------- 选项C ----------
		      \item \textbf{C. <\%@ \%> — JSP指令}
		      \begin{itemize}[label={}, leftmargin=*, nosep]
			      \item \textbf{类型}：Directive（指令）
			      \item \textbf{作用}：页面设置、引入文件等
			            \begin{lstlisting}[language=Java]
<%@ page contentType="text/html;charset=UTF-8" %>
<%@ include file="header.jsp" %>
<%@ taglib prefix="c" uri="http://java.sun.com/jsp/jstl/core" %>
      \end{lstlisting}
		      \end{itemize}

		      % ---------- 选项D ----------
		      \item \textbf{D. <jsp:param> — JSP动作元素}
		      \begin{itemize}[label={}, leftmargin=*, nosep]
			      \item \textbf{类型}：JSP Action
			      \item \textbf{作用}：传递参数
			            \begin{lstlisting}[language=Java]
<jsp:include page="header.jsp">
    <jsp:param name="title" value="首页" />
</jsp:include>
      \end{lstlisting}
		      \end{itemize}

		      \end{itemize}

					  % ======== 阶段三：应背诵知识点 ========
		      \textbf{【应背诵知识点（命题人思维版）】}
		      \begin{enumerate}[leftmargin=*, nosep]
			      \item \textbf{JSP语法分类}：
			            \begin{tabular}{|c|c|c|}
				            \hline
				            \textbf{类型} & \textbf{语法} & \textbf{示例}             \\
				            \hline
				            指令          & <\%@ \%>    & page, include, taglib   \\
				            \hline
				            动作          & <jsp:>      & include, forward, param \\
				            \hline
				            脚本          & <\% \%>     & 代码片段                    \\
				            \hline
				            表达式         & <\%= \%>    & 输出                      \\
				            \hline
				            声明          & <\%! \%>    & 方法定义                    \\
				            \hline
			            \end{tabular}

			      \item \textbf{高频陷阱清单}：
			            \begin{itemize}[label={}, nosep]
				            \item 陷阱1：<jsp:include>是动作，不是指令
				            \item 陷阱2：指令只有三种：page, include, taglib
			            \end{itemize}
		      \end{enumerate}

		      % ======== 阶段四：命题趋势与实战策略 ========
		      \textbf{【命题趋势与实战策略】}

		      \textbf{考场秒杀技巧}：指令用<\%@ \%>，动作用<jsp:>
	      \end{bluebox}

\end{enumerate}

\end{document}
